\section{The Claim}

Most accounts of the world begin with things, particles or fields or bits, and then try to explain how rich structure is built from them. This account begins one level lower, with a single kind of difference, and shows that a long and specific list of known physics is forced once you take that difference seriously and demand it combine without ever losing information.

The single substance carries one property, a value that is one of three: a low value, a middle value, and a high value, written minus one, zero, and plus one. Nothing else is put in at the base. Not space, not time, not particles, not forces. There is only this three-valued stuff, arranged on a growing mesh and shuffled by one fixed local rule. Everything the rest of the paper builds, spacetime and matter and gravity and the quantum, is a large-scale pattern in that one growing arrangement. Where one slogan of physics is "it from bit," the world built from yes-or-no answers, this is "it from difference," the world built from a single distinction taken seriously.

The model rests on two axioms. The first names the substance, the second names its motion.

\begin{axiom}[The vibe]
The one substance is the vibe, which is experience. A vibe's tone is the quality of that experience, not a thing apart from it. There is nothing beneath it.
\end{axiom}

\begin{axiom}[The flow]
The whole is the flow, the eight base elements running together as one, the vibe in endless motion.
\end{axiom}

The first axiom is the experiential posit, and it is the one place the model reaches past physics. Nothing in the derivations below depends on the substance being felt rather than merely a ternary value, so a reader who wants only the physics and the computation may read the vibe as a plain three-valued charge and lose nothing. The point of stating it is that the same structure which yields the physics is, on the fuller reading, the structure of experience itself.

The base is exactly eight elements, no more. Three give the where, a space and its parts. Two give the what, a value and its order. Two give the how, one law for the state and one for the space. The last gives the when. Everything in the paper is a pattern made of these eight, and nothing is ever added underneath them.

\begin{table}[ht]
\centering
\small
\begin{tabularx}{\textwidth}{@{}llX@{}}
\toprule
element & role & what it is \\
\midrule
\texttt{mesh} & where & the whole space, the growing $\{3,4,3,4\}$ hyperbolic honeycomb \\
\texttt{dock} & where & one cell of the mesh, a here-now where tones meet and move on \\
\texttt{site} & where & one of a dock's twenty-four directions, a $D_4$ root, holding one tone \\
\texttt{tone} & what & the value at a site, one of $\{-1, 0, +1\}$ \\
\texttt{lean} & what & the order $-1 < 0 < +1$, turning a tone into signed charge and giving growth its direction \\
\texttt{knit} & how & the law of the state, collide then stream, once a beat \\
\texttt{wake} & how & the law of the space, new docks born at the growing edge \\
\texttt{beat} & when & one discrete tick of time \\
\bottomrule
\end{tabularx}
\caption{The eight base elements. The whole model is these and nothing else.}
\label{tab:elements}
\end{table}

Three claims are made, in order of strength. First, that a single law on a single substrate can be written down with no hidden parts, no tunable dials, nothing imported from known physics, established simply by exhibiting it. Second, that running this law reproduces a specific list of known physics, a list long enough and sharp enough that it could have come out wrong: the exact fractional charges of the quarks, the gauge group of the Standard Model, the weak mixing angle, the spin-statistics connection, the electron's magnetic moment, the inverse-square law of gravity, the bending of light by the general-relativistic amount, the Tsirelson bound of quantum correlations. Third, that the same mesh is a universal computer with a solved address for every cell, so the physical and the computational substrate are one object read two ways.

The work builds on a tradition. The connection between octonions and quark color goes back to G\"unaydin and G\"ursey \cite{gunaydin1973}, the modern division-algebra construction of one generation is due to Furey \cite{furey2018}, the warped-hierarchy mechanism is Randall and Sundrum \cite{randall1999}, the reading of the Higgs as an internal connection is from Connes \cite{connes1996}, the octonion and triality background follows Baez \cite{baez2002octonions}, and the discrete hyperbolic substrate, its cell-addressing, and the whole idea of computing on a hyperbolic honeycomb draw heavily on Margenstern's two volumes on cellular automata in hyperbolic spaces \cite{margenstern2007vol1, margenstern2008vol2}. The contribution here is to chain these into one forced line from a single premise, to add the discrete reversible law and the growing mesh, and to check every link by direct computation in a reproducible suite \cite{cluesurf2026repo}. The broader framework in which this physical core sits, including its account of life, mind, and experience, is the companion monograph \cite{pollard2026vibe}.

This is exploratory, feasibility-stage work. It is offered as a direction worth pursuing, not as a finished theory, and the chapters that follow are as careful to mark what fails and what is open as they are to claim what works.
